\section{Discussion and Future Works}
\label{sec:discussion}

\ourslong (\ours) introduces a robust method leveraging neural fields to eliminate feature artifacts from ViTs. This work additionally pinpoint positional embeddings as the primary source of these artifacts, despite their importance in various vision tasks. Using a neural field optimization process, \ours efficiently extracts clean features from the noise-riddled feature maps of existing ViTs. And using a scalable feature denoiser model, \ours eliminates the need for individual image optimizations. When learned from individually denoised samples, our denoiser generalizes well to unseen data and improves pre-trained ViTs by large margins in dense vision tasks. More broadly, our research suggests several avenues for future exploration: (1) understanding the role of positional embeddings in ViT could inform the design of next-generation deep learning architectures, and (2) redefining positional embeddings within ViTs and transformers is also an imperative problem. Lastly, combining the insights from our work and those of \cite{darcet2023vision} could lead to a \textit{more complete} picture of how these artifacts emerge. We hope that the results presented in this work contribute to a deeper understanding of artifacts in vision transformers and beyond.
